%! AP Statistics — Unit 2, Lesson 2.1 — Homework
\documentclass[10pt]{article}
\usepackage{apstats-article}
\usepackage{apstats-boxes}

%=============================================================================
\begin{document}

\pageheader{Unit 2, Lesson 2.1}{Homework}
\namedateperiod

\vspace{0.1in}

\begin{tcolorbox}[colback=sky, colframe=navy, boxrule=0.8pt, arc=2mm,
    left=3mm, right=3mm, top=1.5mm, bottom=1.5mm]
\small For each news blurb: (a) identify the two variables and their types,
(b) evaluate whether the apparent association is likely \textbf{real} or possibly
\textbf{random}, and (c) identify one \textbf{confounding variable}.
Write in complete phrases.
\end{tcolorbox}

\vspace{0.14in}

% ════════════════════════════════════════════════════════════════════════════
% BLURB 1
% ════════════════════════════════════════════════════════════════════════════
\begin{blurbbox}[Study 1]{navylight}
\begin{headlinebox}{sky}
\small\textit{``A national study of 4{,}000 adults found that people who reported
exercising at least 3 times per week had significantly lower rates of depression
than those who exercised less.''}
\end{headlinebox}

\vspace{8pt}
\renewcommand{\arraystretch}{1.8}
\begin{tabularx}{\linewidth}{@{} l X @{}}
Two variables \& types: & \blank{9cm} \\
Real or possibly random? & \blank{9cm} \\
One confounding variable: & \blank{9cm} \\
\end{tabularx}
\end{blurbbox}

\vspace{0.14in}

% ════════════════════════════════════════════════════════════════════════════
% BLURB 2
% ════════════════════════════════════════════════════════════════════════════
\begin{blurbbox}[Study 2]{greenacc}
\begin{headlinebox}{greenbg}
\small\textit{``In a sample of 15 cities, researchers noticed that cities with more
libraries also tended to have lower crime rates.''}
\end{headlinebox}

\vspace{8pt}
\renewcommand{\arraystretch}{1.8}
\begin{tabularx}{\linewidth}{@{} l X @{}}
Two variables \& types: & \blank{9cm} \\
Real or possibly random? & \blank{9cm} \\
One confounding variable: & \blank{9cm} \\
\end{tabularx}
\end{blurbbox}

\vspace{0.14in}

% ════════════════════════════════════════════════════════════════════════════
% BLURB 3
% ════════════════════════════════════════════════════════════════════════════
\begin{blurbbox}[Study 3]{redacc}
\begin{headlinebox}{redbg}
\small\textit{``Researchers analyzed data from 10{,}000 people and found that
individuals who sleep fewer than 6 hours per night are more likely to be overweight
than those who sleep 7--9 hours.''}
\end{headlinebox}

\vspace{8pt}
\renewcommand{\arraystretch}{1.8}
\begin{tabularx}{\linewidth}{@{} l X @{}}
Two variables \& types: & \blank{9cm} \\
Real or possibly random? & \blank{9cm} \\
One confounding variable: & \blank{9cm} \\
\end{tabularx}
\end{blurbbox}

\vspace{0.14in}

% ── REFLECTION ───────────────────────────────────────────────────────────────
\begin{reflectionbox}
\small Choose \textbf{one} of the three studies above and answer:
\textit{Why isn't the observed association enough to prove causation?}
What additional evidence would you need before claiming one variable causes the other?

\vspace{6pt}
\textcolor{slate}{\scriptsize Study \#:\;\blank{0.6cm}}

\vspace{4pt}
\writeline\\[1pt]\writeline\\[1pt]\writeline\\[1pt]\writeline
\end{reflectionbox}

\vspace{0.14in}

% ── EXTENSION ────────────────────────────────────────────────────────────────
\begin{extensionbox}
\small

\textbf{Spurious Correlations}\\[4pt]
The website \textit{tylervigen.com/spurious-correlations} shows pairs of variables
that are highly correlated in real data but clearly have no causal relationship.

\vspace{5pt}
\begin{enumerate}[leftmargin=1.5em, itemsep=8pt]
  \item Describe one spurious correlation from the website (or make up a plausible one).
        What are the two variables?\\[4pt]
        \writeline\\[1pt]\writeline

  \item Explain why this association almost certainly does not reflect a real relationship.
        What might explain it (e.g., both variables trend over time)?\\[4pt]
        \writeline\\[1pt]\writeline

  \item \textbf{Preview:} In Lesson 2.2 we analyze two \emph{categorical} variables using
        a two-way table. Think of a two-categorical-variable research question you are curious about.\\[4pt]
        \writeline\\[1pt]\writeline
\end{enumerate}

\end{extensionbox}

\end{document}
